\begin{pastexam}[label={px:ntut-115-ee-ds-20}]{北科大電機所（戊組）115 學年度 簡答題第 5 題}
（簡答）

Order the following functions by their growth rate from slowest to fastest: $2^{\sqrt{n}}, n^{1.5},
\log(n!), \frac{n^2}{\log n}, n\log^2 n, 3^n, n^{\log n}$. \textbf{(4\%)}
\end{pastexam}

\begin{pastsoln}
答案（由慢到快）：
\[
\log(n!) \;<\; n \log^2 n \;<\; n^{1.5} \;<\; \frac{n^2}{\log n} \;<\; n^{\log n} \;<\; 2^{\sqrt n} \;<\; 3^n .
\]

先把每個函數化成好比較的形式（資結書第~1.6~節）：
\begin{itemize}
\item $\log(n!) = \Theta(n \log n)$：$n! \le n^n$ 得 $\log(n!) \le n \log n$；$n! \ge (n/2)^{n/2}$ 得
$\log(n!) \ge \frac n2 \log \frac n2$。
\item 三個多項式等級的比：$n \log n < n \log^2 n$（多一個 $\log n$）；$n \log^2 n < n^{1.5}$，因為
$\log^2 n = o(n^{0.5})$（任何 $\log$ 的次方都比 $n$ 的任何正次方慢）；$n^{1.5} < n^2/\log n$，因為兩者相除是
$n^{0.5}/\log n \to \infty$。
\item $n^{\log n} = 2^{\log^2 n}$（以 $2$ 為底）：比任何多項式快（指數 $\log^2 n$ 會超過任何常數 $c \log n$），但比
$2^{\sqrt n}$ 慢，因為 $\log^2 n = o(\sqrt n)$。 例如 $n = 2^{64}$ 時 $\log^2 n = 4096$，$\sqrt n = 2^{32}$。
\item $2^{\sqrt n} < 3^n$：取 $\log_2$ 比較，$\sqrt n$ 對 $n \log_2 3$，後者快得多。
\end{itemize}
注意 $n^{\log n}$ 和 $2^{\sqrt n}$ 在 $n$ 小的時候大小相反（$n = 2^{8}$ 時 $\log^2 n = 64$、$\sqrt n = 16$），成長率談的
是 $n$ 夠大以後（$n > 2^{16}$）的情形。
\end{pastsoln}
